\section{Conclusion}\label{sec:conclusion}
The harmonic coupling of a time-modulated RIS is not unknown interference: it is the product of a
switching timing that the network sets and the receiver knows, and a propagation-plus-dispersion
part that it does not. Grouping the harmonic expansion by switching state makes this separation
exact. Built on it, a receiver needs no element model; estimates the coupling from two to eight pilot
symbols in a low-dimensional shared delay subspace; finds the RIS clock offset by itself, once one
recognizes that the offset is an opposite delay of the two hops and uses the evidence rather than
the fit; learns the static base-station side once; updates cheaply; and equalizes without forming the
coupling matrix. Against the strongest baselines we could build, including an ideal-element estimator with our prior and robust
equalizers for all, it is ahead on every channel and pilot count, exceeds with two pilot symbols a
full-harmonic basis expansion given 32,
decodes where the conventional receiver never does, and after paying for its pilots keeps
3.8--7.4\,bit/s/Hz (held-out channels) on 3GPP channels from walking speed to 120\,km/h. The limits found, delay spreads
beyond the cyclic prefix, high-speed users at high SNR, and the need to know every panel's pattern,
point to the natural next steps: a Doppler-aware coupling model, extended-prefix operation, and
over-the-air validation with a real time-modulated surface.
